\subsection{The \texttt{fill} Command}

Another common command in \MP{} is the \texttt{fill} command.  This is
used to fill closed paths (or cycles).  In order to construct a cycle,
\texttt{cycle} may be appended to the path declaration.  For example,

\begin{lstlisting}[xleftmargin=7bp]
path p;
p:=z1{right}..z2{dir 45}..{up}z3--cycle;
fill p withcolor red;
draw p;
\end{lstlisting}
produces Figure \ref{fig:fill}.  Notice that \texttt{p} is essentially
the same curved path as in Figure \ref{fig:draw1} with the additional
piece that connects \texttt{z3} back to \texttt{z1} with a line segment
using \texttt{-{}-cycle}.

\begin{figure}
  \centering
  \includegraphics{fill.mps}
  \caption{\texttt{fill} example}
  \label{fig:fill}
  \attach{fill.mp}
\end{figure}

Just as it is necessary to fill closed paths, it may also be necessary
to \textit{unfill} closed paths.  For example, the annulus in Figure
\ref{fig:annulus1} can be constructed by

\begin{lstlisting}[xleftmargin=38bp]
color bbblue;
bbblue:=(3/5,4/5,1);
path p,q;
p:=fullcircle scaled (2*54);
q:=fullcircle scaled (2*27);
fill p withcolor bbblue;
unfill q;
draw p;
draw q;
\end{lstlisting}

The \texttt{fullcircle} path is a built-in path that closely
approximates a circle in \MP{} with diameter 1\,bp traversed
counter-clockwise.  This path is not exactly a circle since it is
parameterized by a B\'{e}zier curve and not by trigonometric functions;
however, visually it is essentially indistinguishable from an exact
circle.

\begin{figure}
  \centering
  \includegraphics{annulus-1.mps}
  \caption{\texttt{unfill} example}
  \label{fig:annulus1}
  \attach{annulus.mp}
\end{figure}

Notice that \texttt{p} is a \texttt{fullcircle} of radius 54\,bp
(3/4\,in) and \texttt{q} is a \texttt{fullcircle} of radius 27\,bp
(3/8\,in).  The annulus is constructed by filling \texttt{p} with the
baby blue color \texttt{bbblue} and then unfilling \texttt{q}.  The
\texttt{unfill} command above is equivalent to

\begin{center}
  \verb|fill q withcolor background;|
\end{center}
where \texttt{background} is a built-in color which is \texttt{white} by
default.

Often the \texttt{unfill} command appears to be the natural method for
constructing figures like Figure \ref{fig:annulus1}.  However, the
\texttt{fill} and \texttt{unfill} commands in Figure \ref{fig:annulus1}
can be replaced by

\begin{center}
  \verb|fill p--reverse q--cycle withcolor bbblue;|
\end{center}

\begin{figure}
  \centering
  \includegraphics{annulus-2.mps}
  \caption{Avoiding an \texttt{unfill}}
  \label{fig:annulus2}
  \attach{annulus.mp}
\end{figure}

The path \verb|p--reverse q--cycle| travels around \texttt{p} in a
counter-clockwise directions (since this is the direction that
\texttt{p} traverses) followed by a line segment to connect to
\texttt{q}.  It then traverses clockwise around \texttt{q} (using the
\texttt{reverse} operator) and finally returns to the starting point
along a line segment using \texttt{-{}-cycle}.  This path is illustrated
in Figure~\ref{fig:annulus2}.  One reason for using this method to
construct the annulus as opposed to the \texttt{unfill} command is to
ensure \textit{proper transparency} when placing the figure in an
external document with a non-white background.  If the former method is
used and the annulus is placed on a non-white background, say magenta,
then the result is Figure \ref{fig:annulus3}.

\begin{figure}
  \centering
  \includegraphics{annulus-3.mps}
  \caption{Improper transparency using \texttt{unfill}}
  \label{fig:annulus3}
  \attach{annulus.mp}
\end{figure}

It may be desired to have the interior of \texttt{q} be magenta instead
of \texttt{white}.  This could be accomplished by redefining
\texttt{background}; however, the latter method described above is a
much simpler solution.
